\documentclass[a4paper, 11pt]{article}
\usepackage{xeCJK}
\usepackage{geometry}
\usepackage{titlesec}

\geometry{margin=1in}
\setCJKmainfont{MS Mincho}

% \title{進捗報告}
% \author{コレドール・パブロ}
% \date{２０２５年１０月１０日}

\titlespacing*{\section}{0pt}{2ex}{1ex}
\titlespacing*{\subsection}{0pt}{1ex}{0.8ex}
\titlespacing*{\subsubsection}{0pt}{1ex}{0.5ex}

\begin{document}

\begin{center}
    {\LARGE \bfseries 進捗報告} \\ % Title: Large and bold
    \vspace{0.3em} % A small vertical space
    Juan Pablo Corredor \\ % Author
    ２０２５年１０月２４日 % Date
\end{center}
\vspace{0.5em} % Adjust space between title block and main text

\section*{近況}
% Write your recent findings here.
\begin{itemize}
    \item カタロニア映像大学のBGM制作が進んでいる。
    \item 藝大のアニメ科のBGMも進んでいる。
    \item 神天気で嬉しい。
\end{itemize} 

\section*{やったこと}
\begin{itemize}
    \item コードに集中し、まだ開発中。
    \item Inharmonicityという概念について調べてみた。\cite{houtsma1982inharmonicity} \newline \cite{murray2022inharmonicity}も読んでみたいい。
\end{itemize}


\section*{考えていること}

トレーニングデータを生成するプログラムの構造を大分改良し、現在状況は以下の通り：

\begin{itemize}
    \item 音程認識 \newline Pyinアルゴリズムを使用。分析しない可能性もあるため、その場合はファイルのメタデータに含まれている音程が使用される。
    \item 周期認識 \newline 相互相関分析で各周期のゼロクロスサンプルを探り、行列に記録。分析しない場合はFFTのHop Sizeの初サンプルが記録される。
    \item 代表的な倍音を探る \newline 音源のトランジェント数ミリ秒後から長い（65536サンプル）FFTによって最も代表的な倍音を探る。分析しない場合は32倍音未満は順番に記録される。
    \item 倍音エンベロープ分析 \newline 周期リスト、またはHop Sizeを従い、FFT分析で各倍音のエンベロープが生成される。FFTウィンドウには3つの異なる方法が使用されている：
    \begin{itemize}
        \item Audio Chunk: 音源からFFTウィンドウにデータを直接に書き込む。
        \item Single Period: 1周期のみFFTウィンドウに書き込む。ウィンドウの残りはゼロパッドされる。
        \item Period Loop: FFTウィンドウが埋まるまで、1周期の波形を繰り返す。
    \end{itemize}
    \item ノイズ分析 \newline 方法検討中
    \item トランジェント分析 \newline 方法検討中
\end{itemize}

トレーニングに向かうなお更、モデルの構造を考えている。\cite{latentSpaceInterpolation}を参考にし、シンセサイザーの「パラメーター」を倍音エンベロープとして、生成できるモデルを開発する。そのため、倍音エンベロープに留まらず、スペクトルグラムも使用する予定。


\section*{やること}
% Outline future steps here.
\begin{itemize}
    \item 開発を続ける。
    \item ノイズ分析の論文を探す。
\end{itemize}

% Add your bibliography here.
\bibliography{../../ref}
\bibliographystyle{apalike}

\end{document}